\section{The Jacobian}

Our goal is to decompose \(\Jac(X)\) up to \(\QQ\)-isogeny by exploiting the \(S_4\)-module structure on \(H^0(X,\Omega_X^1)\). We first work out this structure and then identify each isotypic factor with a concrete abelian variety.

\subsection{Differentials}

\begin{theorem}\label{thm:hodge}
Let \(V_2\) be the two-dimensional irreducible representation of \(S_4\), \(V_3\) the standard three-dimensional representation, and \(V_3^-:=V_3\otimes\mathrm{sgn}\).
Then
\[
H^0(X,\Omega_X^1)\cong \mathrm{sgn}^{\oplus 2}\oplus V_2\oplus V_3\oplus (V_3^-)^{\oplus 3}.
\]
Accordingly there are rational isotypic factors \(A_{\mathrm{sgn}}, A_{V_2}, A_{V_3}\), and \(A_{V_3^-}\) such that
\[
\Jac(X)\sim_{\QQ}A_{\mathrm{sgn}}\times A_{V_2}\times A_{V_3}\times A_{V_3^-},
\]
with dimensions
\[
\dim A_{\mathrm{sgn}}=2,\qquad \dim A_{V_2}=2,\qquad \dim A_{V_3}=3,\qquad \dim A_{V_3^-}=9.
\]
\end{theorem}

\begin{proof}
Since \(X/S_4\cong \PP^1\), there is no trivial summand. Write
\[
H^0(X,\Omega_X^1)\cong m_{\mathrm{sgn}}\cdot \mathrm{sgn}\oplus m_2\cdot V_2\oplus m_3\cdot V_3\oplus m_3^-\cdot V_3^-.
\]
For any subgroup \(H\subset S_4\), the \(H\)-invariants have dimension \(g(X/H)\).

From Theorem~\ref{thm:quotients}, \(g(X/A_4)=2\) and \(g(X/V_4)=4\). If \(S_3\subset S_4\) is a point stabilizer, then \(X/S_3\cong E\), so \(g(X/S_3)=1\). Let \(C=\langle(1234)\rangle\) and \(D=N_{S_4}(C)\cong D_4\). Character averages give
\[
\begin{array}{c|ccccc}
 & A_4 & V_4 & S_3 & C & D\\ \hline
\mathrm{sgn} & 1 & 1 & 0 & 0 & 0\\
V_2 & 0 & 2 & 0 & 1 & 1\\
V_3 & 0 & 0 & 1 & 0 & 0\\
V_3^- & 0 & 0 & 0 & 1 & 0
\end{array}
\]
for the invariant dimensions of each irreducible constituent.

Reading off the \(A_4\)- and \(V_4\)-columns gives
\[
m_{\mathrm{sgn}}=2,\qquad m_{\mathrm{sgn}}+2m_2=4,
\]
hence \(m_2=1\), and the \(S_3\)-column gives \(m_3=1\). The genus constraint \(g(X)=16\) then forces
\[
16=m_{\mathrm{sgn}}+2m_2+3m_3+3m_3^-,
\]
so \(m_3^-=3\). The Jacobian decomposition follows from the rational isotypic decomposition.
\end{proof}

\subsection{The dihedral quotient}

\begin{proposition}\label{prop:eres}
Let \(C=\langle(1234)\rangle\subset S_4\), let \(D=N_{S_4}(C)\cong D_4\), and set
\[
Y:=X/C,\qquad E_{\mathrm{res}}:=X/D.
\]
Then:
\begin{enumerate}
\item[(i)] The function field of \(E_{\mathrm{res}}\) is
\[
\QQ(E_{\mathrm{res}})=\QQ(y,z),\qquad R(z,y)=0,
\]
where \(R\) is the cubic resolvent \eqref{eq:resolvent}.
\item[(ii)] The smooth projective model of \(R(z,y)=0\) is birational over \(\QQ\) to the elliptic curve
\[
 \mathcal E_{\mathrm{res}}:\quad w_{\mathrm{res}}^2=x_{\mathrm{res}}^3-16x_{\mathrm{res}}^2-64x_{\mathrm{res}}+1040.
\]
In particular,
\[
g(E_{\mathrm{res}})=1,\qquad
\Delta(\mathcal E_{\mathrm{res}})=2^{12}\cdot 31^2\cdot 37,\qquad
j(\mathcal E_{\mathrm{res}})=\frac{2^{18}\cdot 7^3}{31^2\cdot 37}.
\]
\item[(iii)] \(g(Y)=4\), and the natural map \(Y\to E_{\mathrm{res}}\) has degree \(2\) and is branched at exactly six points.
\end{enumerate}
\end{proposition}

\begin{proof}
Since \(V_4\triangleleft D\) and \(D/V_4\) has order \(2\), the quotient \(E_{\mathrm{res}}=X/D\) sits inside the resolvent field \(\QQ(X_V)/\QQ(y)\). Under the identification \(S_4/V_4\cong S_3\), the image of \(D\) stabilizes one of the three pairings of the four quartic roots, so the fixed field is generated by a single resolvent root \(z\):
\[
\QQ(E_{\mathrm{res}})=\QQ(y,z),\qquad R(z,y)=0.
\]

The affine model \(R(z,y)=0\) passes through \((y,z)=(0,1)\). Projecting from this point, set
\[
u=\frac{z-1}{y}.
\]
Substituting \(z=1+uy\) into \(R(z,y)=0\) gives
\[
(u^3+2u^2-4u-8)y^2+(4u^2-17)y+(4u-7)=0.
\]
Set
\[
v=2(u^3+2u^2-4u-8)y+(4u^2-17).
\]
Then the quadratic discriminant identity becomes
\[
v^2=-4u^3-16u^2+16u+65.
\]
After the change of variables
\[
 x_{\mathrm{res}}=-4u,\qquad w_{\mathrm{res}}=4v,
\]
we obtain the Weierstrass model
\[
 w_{\mathrm{res}}^2=x_{\mathrm{res}}^3-16x_{\mathrm{res}}^2-64x_{\mathrm{res}}+1040.
\]
The inverse rational map is recovered from
\[
 u=-\frac{x_{\mathrm{res}}}{4},\qquad
 v=\frac{w_{\mathrm{res}}}{4},\qquad
 y=\frac{v-(4u^2-17)}{2(u^3+2u^2-4u-8)},\qquad
 z=1+uy.
\]
The normalization of \(R(z,y)=0\) is therefore an elliptic curve, and the discriminant and \(j\)-invariant are read off from the integral model.

From the character table in Theorem~\ref{thm:hodge}, the invariant dimensions for \(C\) and \(D\) are
\[
g(Y)=\dim H^0(X,\Omega_X^1)^C=1+3=4,\qquad
g(E_{\mathrm{res}})=\dim H^0(X,\Omega_X^1)^D=1.
\]
Since \([D:C]=2\), the map \(Y\to E_{\mathrm{res}}\) has degree \(2\). Hurwitz gives
\[
2g(Y)-2=2(2g(E_{\mathrm{res}})-2)+r=r,
\]
so the total ramification degree is \(r=6\).
\end{proof}

\begin{proposition}\label{prop:eres-conductor}
The minimal discriminant and conductor of \(\mathcal E_{\mathrm{res}}\) are
\[
\Delta_{\min}(\mathcal E_{\mathrm{res}})=31^2\cdot 37,
\qquad
N(\mathcal E_{\mathrm{res}})=31\cdot 37=1147.
\]
In particular, \(\mathcal E_{\mathrm{res}}\) has multiplicative reduction at both \(31\) and \(37\).
\end{proposition}

\begin{proof}
A short computation gives \(c_4=2^{10}\cdot 7\), \(c_6=-2^9\cdot 23\cdot 29\), and \(\Delta=2^{12}\cdot 31^2\cdot 37\).
Since \(v_2(c_4)=10\ge 4\), \(v_2(c_6)=9\ge 6\), and \(v_2(\Delta)=12\ge 12\), one Tate reduction step with \(u=2\) yields the minimal model with
\[
c_4'=c_4/2^4=448=2^6\cdot 7,
\qquad
c_6'=c_6/2^6=-5336=-2^3\cdot 23\cdot 29,
\qquad
\Delta'=\Delta/2^{12}=31^2\cdot 37.
\]
Since \(v_2(c_6')=3<6\), no further reduction at~\(2\) is possible; at all odd primes, \(v_p(c_4')<4\) or \(v_p(\Delta')=0\), so the model is globally minimal.
Now \(v_{31}(c_4')=0\) and \(v_{31}(\Delta')=2>0\), so the reduction at \(31\) is multiplicative (conductor exponent~\(1\)).
Similarly \(v_{37}(c_4')=0\) and \(v_{37}(\Delta')=1>0\), giving multiplicative reduction at \(37\).
\end{proof}

\begin{proposition}\label{prop:extract}
Let \(G\) be a finite group acting on a smooth projective curve \(X\) over \(\QQ\), and let \(W\) be an absolutely irreducible rational representation of \(G\) such that \(\End_G(W)=\QQ\). Write the \(W\)-isotypic component of \(H^0(X,\Omega_X^1)\) as
\[
H^0(X,\Omega_X^1)_W\cong W\otimes M_W,
\]
where \(M_W\) is the multiplicity space.
If \(A_W\subset \Jac(X)\) denotes the corresponding rational isotypic factor, then there exists an abelian variety \(B_W\), unique up to \(\QQ\)-isogeny, such that
\[
H^0(B_W,\Omega_{B_W}^1)\cong M_W,
\qquad
A_W\sim_{\QQ} B_W^{\dim W}.
\]
In particular, whenever \(M_W\) is realized by the holomorphic differentials of a quotient or Prym variety \(B\), one has
\[
A_W\sim_{\QQ} B^{\dim W}.
\]
\end{proposition}

\begin{proof}
We recall the standard group-algebra argument underlying \cite{KaniRosen1989}; see also \cite[Section~13.4]{BirkenhakeLange2004}.

Let \(e_W\in \QQ[G]\) be the central idempotent attached to \(W\). Since \(\End_G(W)=\QQ\), the corresponding Wedderburn block is
\[
\QQ[G]e_W\cong M_d(\QQ),\qquad d=\dim W.
\]
Choose a primitive idempotent \(\varepsilon\in M_d(\QQ)\), and define
\[
B_W:=\varepsilon\cdot A_W.
\]
The central idempotent \(e_W\) is a sum of \(d\) pairwise orthogonal conjugates of \(\varepsilon\), so
\[
A_W\sim_{\QQ} B_W^d.
\]
On differentials, \(H^0(A_W,\Omega_{A_W}^1)\cong W\otimes M_W\), and \(\varepsilon\) picks out a single matrix-column factor:
\[
H^0(B_W,\Omega_{B_W}^1)\cong M_W.
\]
Any other abelian variety \(B\) realizing the same multiplicity space defines the same rational Hodge substructure of \(H^1(X)\), hence is \(\QQ\)-isogenous to~\(B_W\).
\end{proof}

\begin{remark}\label{rem:schur}
Every complex irreducible representation of \(S_4\) is realizable over \(\QQ\) (Schur index \(1\)), so the hypothesis \(\End_G(W)=\QQ\) of Proposition~\ref{prop:extract} is automatically fulfilled for each of the four isotypic constituents.
\end{remark}

\begin{proposition}\label{prop:sgn-v2}
Let \(W_{\mathrm{sgn}}\) and \(W_{V_2}\) be the sign and \(V_2\)-isotypic summands of \(H^0(X,\Omega_X^1)\). Then
\[
H^0(X,\Omega_X^1)^{A_4}=W_{\mathrm{sgn}},
\qquad
H^0(X,\Omega_X^1)^{V_4}=W_{\mathrm{sgn}}\oplus W_{V_2}.
\]
Moreover, \(V_4\) acts trivially on \(V_2\), and the quotient group \(A_4/V_4\cong C_3\) acts trivially on \(W_{\mathrm{sgn}}\) and through the unique nontrivial irreducible rational representation on \(W_{V_2}\). Consequently,
\[
A_{\mathrm{sgn}}\sim_{\QQ}\Jac(X_A),
\qquad
H^0(P,\Omega_P^1)\cong W_{V_2},
\qquad
P\sim_{\QQ}A_{V_2}.
\]
\end{proposition}

\begin{proof}
From the character table computed in Theorem~\ref{thm:hodge}, the only constituent with nonzero \(A_4\)-invariants is \(\mathrm{sgn}\), contributing one invariant line per copy. Hence
\[
H^0(X,\Omega_X^1)^{A_4}=W_{\mathrm{sgn}}.
\]
For \(V_4\), the same table shows that the only constituents with nonzero invariants are \(\mathrm{sgn}\) and \(V_2\), with invariant dimensions \(1\) and \(2\), respectively. Therefore
\[
H^0(X,\Omega_X^1)^{V_4}=W_{\mathrm{sgn}}\oplus W_{V_2}.
\]

The character of \(V_2\) evaluates to \(2\) on every element of \(V_4\), so \(V_4\) acts trivially on \(W_{V_2}\). Since \(V_2\) has no \(A_4\)-invariants, the quotient \(A_4/V_4\cong C_3\) acts without fixed vectors, forcing \(W_{V_2}\) to be the unique nontrivial irreducible rational \(C_3\)-module. The sign representation is trivial on \(A_4\), so the action on \(W_{\mathrm{sgn}}\) is trivial.

Pullback identifies \(H^0(X_A,\Omega_{X_A}^1)\) with \(H^0(X,\Omega_X^1)^{A_4}=W_{\mathrm{sgn}}\). Since \(\mathrm{sgn}\) is one-dimensional, Proposition~\ref{prop:extract} gives
\[
A_{\mathrm{sgn}}\sim_{\QQ}\Jac(X_A).
\]
Likewise,
\[
H^0(X_V,\Omega_{X_V}^1)=H^0(X,\Omega_X^1)^{V_4}=W_{\mathrm{sgn}}\oplus W_{V_2}.
\]
By Corollary~\ref{cor:prym},
\[
H^0(X_V,\Omega_{X_V}^1)\cong \pi^*H^0(X_A,\Omega_{X_A}^1)\oplus H^0(P,\Omega_P^1),
\]
where \(\pi:X_V\to X_A\) is the natural cover. Comparing the two decompositions shows that
\[
H^0(P,\Omega_P^1)\cong W_{V_2}.
\]
It follows that \(P\) and \(A_{V_2}\) define the same rational Hodge substructure of \(H^1(X)\), so \(P\sim_{\QQ}A_{V_2}\).
\end{proof}

\subsection{The factorization}

\begin{theorem}\label{thm:full-factor}
Let \(E=X/S_3\) be the conductor-\(37\) elliptic quotient from Theorem~\ref{thm:elliptic}, let \(E_{\mathrm{res}}=X/D_4\) be the resolvent elliptic quotient from Proposition~\ref{prop:eres}, and let
\[
Q:=\Prym(Y/E_{\mathrm{res}}),\qquad Y=X/C_4.
\]
Then one has
\[
A_{\mathrm{sgn}}\sim_{\QQ}\Jac(X_A),\qquad
P\sim_{\QQ} E_{\mathrm{res}}^2,\qquad
A_{V_3}\sim_{\QQ} E^3,\qquad
A_{V_3^-}\sim_{\QQ} Q^3.
\]
Consequently,
\[
\Jac(X)\sim_{\QQ}\Jac(X_A)\times E_{\mathrm{res}}^2\times E^3\times Q^3.
\]
\end{theorem}

\begin{proof}
Proposition~\ref{prop:sgn-v2} identifies the sign constituent:
\[
A_{\mathrm{sgn}}\sim_{\QQ}\Jac(X_A).
\]

We identify the remaining three factors in turn.

For \(V_3\), let \(H\subset S_4\) be a point stabilizer, so that \(X/H\cong E\). The invariant table shows that \(V_3\) is the unique constituent with nonzero \(H\)-invariants, giving
\[
H^0(E,\Omega_E^1)\cong M_{V_3},
\]
where \(M_{V_3}\) is the multiplicity space of \(V_3\). Proposition~\ref{prop:extract} gives
\[
A_{V_3}\sim_{\QQ} E^3.
\]

For \(V_2\), we already know \(P\sim_{\QQ}A_{V_2}\) from Proposition~\ref{prop:sgn-v2}. The quotient \(E_{\mathrm{res}}=X/D\) satisfies
\[
H^0(E_{\mathrm{res}},\Omega_{E_{\mathrm{res}}}^1)=H^0(X,\Omega_X^1)^D,
\]
and the \(D\)-invariant column shows that this is the multiplicity space \(M_{V_2}\). Proposition~\ref{prop:extract} then gives
\[
A_{V_2}\sim_{\QQ} E_{\mathrm{res}}^2.
\]
Together with Proposition~\ref{prop:sgn-v2}, this yields
\[
P\sim_{\QQ}E_{\mathrm{res}}^2.
\]

The most delicate factor is \(V_3^-\). The degree-\(2\) map \(Y\to E_{\mathrm{res}}\) has a deck involution \(\iota\in D/C\), and the pullback decomposition gives
\[
H^0(Y,\Omega_Y^1)\cong \pi^*H^0(E_{\mathrm{res}},\Omega_{E_{\mathrm{res}}}^1)\oplus H^0(Q,\Omega_Q^1),
\]
where the first summand is the \(\iota\)-invariant part and the second is the \(\iota\)-anti-invariant part.
Now
\[
H^0(Y,\Omega_Y^1)=H^0(X,\Omega_X^1)^C.
\]
The character table gives \(\dim H^0(X,\Omega_X^1)^C=4\) (one line from \(V_2\), one from each of the three copies of \(V_3^-\)) and \(\dim H^0(X,\Omega_X^1)^D=1\) (from \(V_2\) alone). Since \(V_3^-\) has vanishing character average over \(D\), its \(C\)-invariant lines all lie in the \(\iota\)-anti-invariant eigenspace, so
\[
H^0(Q,\Omega_Q^1)\cong M_{V_3^-},
\]
and \(\dim Q=3\). Proposition~\ref{prop:extract} (applicable by Remark~\ref{rem:schur}) now gives
\[
A_{V_3^-}\sim_{\QQ} Q^3.
\]

Assembling the four identifications gives
\[
\Jac(X)\sim_{\QQ}\Jac(X_A)\times E_{\mathrm{res}}^2\times E^3\times Q^3.
\]
\end{proof}

\begin{remark}\label{rem:prym-correspondence}
The isogeny \(P\sim_{\QQ}E_{\mathrm{res}}^2\) is induced by the same rational correspondence on \(X\) that realizes the \(V_2\)-isotypic projector in \(\QQ[S_4]\). On the one hand, if \(\sigma\) generates the deck group of the \'etale cyclic triple cover \(X_V\to X_A\), then the Prym projector \(1-\sigma\) cuts out the anti-invariant part of \(\Jac(X_V)\), namely \(P\). On the other hand, the quotient map \(X\to X/D_4=E_{\mathrm{res}}\) identifies the same multiplicity space with \(H^0(E_{\mathrm{res}},\Omega_{E_{\mathrm{res}}}^1)\). The resulting pullback and norm correspondences induce homomorphisms
\[
P\longrightarrow E_{\mathrm{res}}^2,
\qquad
E_{\mathrm{res}}^2\longrightarrow P
\]
which are inverse up to finite kernel. This is the concrete \(V_2\)-block of the general functorial decomposition formalism for Jacobians of Galois covers \cite{LombardoEtAl2023}. We do not compute the minimal degree of this isogeny here.
\end{remark}

\begin{corollary}\label{cor:l-factor}
There is a corresponding factorization of first cohomology and of Hasse--Weil \(L\)-functions:
\[
H^1(X)\sim H^1(X_A)\oplus H^1(E_{\mathrm{res}})^{\oplus 2}\oplus H^1(E)^{\oplus 3}\oplus H^1(Q)^{\oplus 3},
\]
hence
\[
L(H^1(X),s)=L(H^1(X_A),s)\,L(E_{\mathrm{res}},s)^2\,L(E,s)^3\,L(H^1(Q),s)^3.
\]
\end{corollary}

\begin{proof}
This follows at once from Theorem~\ref{thm:full-factor} by passing to \(\ell\)-adic representations and taking Euler products.
\end{proof}

\begin{proposition}\label{prop:prym-end}
The Prym surface \(P=\Prym(X_V/X_A)\) carries a polarization of type \((1,3)\) and satisfies
\[
\End^0(P)\cong M_2(\QQ),
\qquad
\rho(P_{\overline{\QQ}})=3.
\]
In particular, \(P\) is not of CM type.
\end{proposition}

\begin{proof}
The Prym polarization for an \'etale cyclic triple cover of a genus-\(2\) curve has type \((1,3)\) \cite[Section~12.9]{BirkenhakeLange2004}. Since
\[
P\sim_{\QQ}E_{\mathrm{res}}^2.
\]
by Theorem~\ref{thm:full-factor}, and the \(j\)-invariant \(j(E_{\mathrm{res}})=2^{18}\cdot 7^3/(31^2\cdot 37)\) is not an algebraic integer, the curve \(E_{\mathrm{res}}\) has no complex multiplication \cite[Ch.~II]{Silverman2009}. It follows that
\[
\End^0(P)\cong \End^0(E_{\mathrm{res}}^2)\cong M_2(\QQ).
\]
Isogeny induces an isomorphism on N\'eron--Severi groups after tensoring with~\(\QQ\), so
\[
\rho(P_{\overline{\QQ}})=\rho((E_{\mathrm{res}}^2)_{\overline{\QQ}}).
\]
For the square of a non-CM elliptic curve the geometric Picard number is~\(3\), generated by the two coordinate classes and the diagonal. Since \(\End^0(P)\cong M_2(\QQ)\) is not a commutative semisimple algebra of dimension \(2\dim P\), the surface \(P\) is not of CM type.
\end{proof}

\begin{remark}\label{rem:prym-triple-cover-context}
The cover \(X_V\to X_A\) fits the cyclic triple-cover setting of Lange--Ortega \cite{LangeOrtega2011}. Their analysis of non-cyclic triple coverings of genus-\(2\) curves \cite{LangeOrtegaTriple2011} exhibits a markedly different Prym behavior, with principally polarized Prym surfaces on the non-cyclic side. In the present cyclic example the Prym polarization has type \((1,3)\), and the additional \(S_4\)-symmetry forces the stronger splitting
\[
P\sim_{\QQ}E_{\mathrm{res}}^2.
\]
\end{remark}

\section{Ray classes and weight one}

We turn to the class-field-theoretic description of the splitting field of the branch cubic. Let \(F=\QQ(\beta)\) be the cubic field generated by a root \(\beta\) of \(c(y)\), let \(L\) be the splitting field, and let
\[
K=\QQ(\sqrt{-111})
\]
be the quadratic resolvent field.

\begin{proposition}\label{prop:class-group}
The ideal class group of \(K=\QQ(\sqrt{-111})\) is cyclic of order \(8\):
\[
\Cl(K)\cong C_8.
\]
\end{proposition}

\begin{proof}
Let
\[
\tau:=\frac{1+\sqrt{-111}}{2},
\]
so \(\mathcal O_K=\ZZ[\tau]\) and \(\tau^2-\tau+28=0\). Since \(\Disc(K)=-111\), the primitive reduced positive definite binary quadratic forms of discriminant \(-111\) are
\[
(1,1,28),\ (2,\pm 1,14),\ (3,3,10),\ (4,\pm 1,7),\ (5,\pm 3,6).
\]
There are eight such forms, so \(h_K=8\).

To exhibit a generator, consider the ideal
\[
\mathfrak a:=(2,\tau-1)\subset \mathcal O_K.
\]
A direct ideal multiplication gives
\[
\mathfrak a^2=(4,\tau+3),\qquad
\mathfrak a^4=(16,\tau+11),\qquad
\mathfrak a^8=(1+3\tau).
\]
so \([\mathfrak a]^8=1\) in \(\Cl(K)\). To verify that the order is exactly \(8\), note that the norm form on \(\mathcal O_K\) is
\[
N_{K/\QQ}(x+y\tau)=x^2+xy+28y^2.
\]
If \(\mathfrak a^4\) were principal, it would be generated by an element of norm \(16\). Solving
\[
x^2+xy+28y^2=16
\]
gives only \((x,y)=(\pm 4,0)\), so the only principal ideal of norm \(16\) is \((4)\). But \(\mathfrak a^4=(16,\tau+11)\neq (4)\), so \([\mathfrak a]\) has order exactly~\(8\) and generates \(\Cl(K)\).
\end{proof}

\begin{theorem}\label{thm:conductor}
Let \(\mathfrak p_3\) be the unique prime of \(K\) above \(3\). Then:
\begin{enumerate}
\item[(i)] \(L/K\) is cyclic of degree \(3\).
\item[(ii)] The Artin conductor of \(L/K\) is
\[
\mathfrak f_{L/K}=\mathfrak p_3^2.
\]
\item[(iii)] The relative discriminant is
\[
\mathfrak d_{L/K}=\mathfrak p_3^4.
\]
\item[(iv)] The absolute discriminant of \(L\) is
\[
\Disc(L/\QQ)=3^7\cdot 37^3.
\]
\end{enumerate}
\end{theorem}

\begin{proof}
Since \(c(y)\) is irreducible with nonsquare discriminant (Proposition~\ref{prop:number-field}), the Galois group of \(L/\QQ\) is~\(S_3\), and its unique index-\(2\) subgroup fixes \(K\). Hence \(L/K\) is cyclic of degree~\(3\).

Let \(\rho\) denote the standard two-dimensional Artin representation of \(S_3\). The factorization
\[
\zeta_F(s)=\zeta(s)L(\rho,s)
\]
shows that the Artin conductor of \(\rho\) is
\[
\mathfrak f(\rho)=|\Disc(F/\QQ)|=999.
\]
On the other hand, if \(\chi\) is a nontrivial character of \(\Gal(L/K)\), then
\[
\rho\cong \operatorname{Ind}_{G_K}^{G_\QQ}\chi.
\]
The conductor formula for induction gives \cite[Ch.~VII]{Neukirch1999}
\[
\mathfrak f(\rho)=|\Disc(K/\QQ)|\,N_{K/\QQ}(\mathfrak f_\chi).
\]
Since \(|\Disc(K/\QQ)|=111\), we obtain
\[
N_{K/\QQ}(\mathfrak f_\chi)=\frac{999}{111}=9.
\]
Since \(3\) ramifies in \(K\), there is a unique prime \(\mathfrak p_3\) above \(3\), and \(\mathfrak p_3^2\) is the only ideal of norm \(9\). Therefore
\[
\mathfrak f_\chi=\mathfrak p_3^2.
\]
The two nontrivial characters of a cyclic cubic extension share the same conductor, so
\[
\mathfrak f_{L/K}=\mathfrak p_3^2
\]
and the conductor--discriminant formula \cite[Ch.~VII]{Neukirch1999} gives \(\mathfrak d_{L/K}=\mathfrak p_3^4\). Taking norms,
\[
N_{K/\QQ}(\mathfrak d_{L/K})=N(\mathfrak p_3)^4=3^4.
\]
Therefore
\[
\Disc(L/\QQ)=\Disc(K/\QQ)^3\,N(\mathfrak d_{L/K})
=111^3\cdot 3^4
=3^7\cdot 37^3.
\]
\end{proof}

\begin{remark}\label{rem:disc-distinction}
The two discriminants appearing here refer to different fields and should not be conflated. The cubic field \(F=\QQ(\beta)\) is the non-Galois cubic subfield of the \(S_3\)-extension \(L/\QQ\), and
\[
\Disc(F/\QQ)=-999=-3^3\cdot 37
\]
by Proposition~\ref{prop:number-field}. By contrast, \(L\) is the full splitting field, so its absolute discriminant is obtained from the quadratic resolvent field \(K=\QQ(\sqrt{-111})\) and the relative discriminant of the cubic extension \(L/K\):
\[
\Disc(L/\QQ)=\Disc(K/\QQ)^3\,N_{K/\QQ}(\mathfrak d_{L/K})
=111^3\cdot 3^4
=3^7\cdot 37^3.
\]
Thus items~(ii)--(iv) are mutually consistent: \(\mathfrak f_{L/K}=\mathfrak p_3^2\) implies \(\mathfrak d_{L/K}=\mathfrak p_3^4\), and after taking norms this yields the displayed discriminant of~\(L\), not the discriminant of the cubic subfield~\(F\).
\end{remark}

\begin{theorem}\label{thm:ray-class}
Let \(H_K\) be the Hilbert class field of \(K=\QQ(\sqrt{-111})\), and let \(K_{\mathfrak p_3^2}\) be the ray class field of modulus \(\mathfrak p_3^2\). Then:
\begin{enumerate}
\item[(i)] The ray class group is cyclic:
\[
\Cl_{\mathfrak p_3^2}(K)\cong C_{24}.
\]
\item[(ii)] The ray class field is the compositum
\[
K_{\mathfrak p_3^2}=H_KL,
\qquad
H_K\cap L=K.
\]
\item[(iii)] The extension \(K_{\mathfrak p_3^2}/\QQ\) is Galois with generalized dihedral group
\[
\Gal(K_{\mathfrak p_3^2}/\QQ)\cong C_{24}\rtimes C_2,
\]
where the nontrivial element of \(C_2\) acts on \(C_{24}\) by inversion.
\item[(iv)] The splitting field \(L\) of \(c(y)\) is the unique cubic subextension of \(K_{\mathfrak p_3^2}/K\).
\end{enumerate}
\end{theorem}

\begin{proof}
The standard exact sequence for the modulus \(\mathfrak p_3^2\) reads \cite[Ch.~VI]{Neukirch1999}
\[
1\longrightarrow \mathcal O_K^\times/(\mathcal O_K^\times\cap (1+\mathfrak p_3^2))
\longrightarrow (\mathcal O_K/\mathfrak p_3^2)^\times
\longrightarrow \Cl_{\mathfrak p_3^2}(K)
\longrightarrow \Cl(K)\longrightarrow 1.
\]
Since \(K\) is imaginary quadratic with discriminant \(-111\), the unit group is \(\mathcal O_K^\times=\{\pm 1\}\), and \(-1\not\equiv 1\pmod{\mathfrak p_3^2}\), so its image in \((\mathcal O_K/\mathfrak p_3^2)^\times\) has order \(2\). The residue field of \(\mathfrak p_3\) has size \(3\), so
\[
|(\mathcal O_K/\mathfrak p_3^2)^\times|
=N(\mathfrak p_3^2)\left(1-\frac{1}{N(\mathfrak p_3)}\right)
=9\left(1-\frac13\right)
=6.
\]
Equivalently, because \(3\) is ramified in \(K\), the local kernel is
\[
\frac{(\mathcal O_K/\mathfrak p_3^2)^\times}{\{\pm1\}}
\cong \frac{1+\mathfrak p_3}{1+\mathfrak p_3^2}
\cong \frac{\mathfrak p_3}{\mathfrak p_3^2}
\cong \F_3,
\]
so it is cyclic of order~\(3\). Thus the kernel of \(\Cl_{\mathfrak p_3^2}(K)\to \Cl(K)\) has order \(3\). Since \(\Cl(K)\cong C_8\) by Proposition~\ref{prop:class-group},
\[
|\Cl_{\mathfrak p_3^2}(K)|=3\cdot 8=24.
\]
Let \([\mathfrak a]\) be the generator of \(\Cl(K)\) from Proposition~\ref{prop:class-group}, and let \(u\) be any nontrivial element of the local kernel. The kernel has order coprime to \(8\), so the sequence splits and
\[
\Cl_{\mathfrak p_3^2}(K)\cong C_8\times C_3\cong C_{24},
\]
with generators \([\mathfrak a]\) and \(u\). This proves~(i).

The Hilbert class field satisfies \([H_K:K]=h_K=8\), and \(L/K\) is cyclic cubic of conductor exactly \(\mathfrak p_3^2\) (Theorem~\ref{thm:conductor}). Since \(\gcd(8,3)=1\),
\[
H_K\cap L=K.
\]
Hence
\[
[H_KL:K]=[H_K:K]\,[L:K]=24=[K_{\mathfrak p_3^2}:K],
\]
so \(K_{\mathfrak p_3^2}=H_KL\). This proves~(ii).

The modulus \(\mathfrak p_3^2\) is stable under complex conjugation, so \(K_{\mathfrak p_3^2}/\QQ\) is Galois and conjugation acts on \(\Cl_{\mathfrak p_3^2}(K)\) by inversion \cite[Ch.~9]{Cox2012}, giving~(iii).

For~(iv), \(C_{24}\) has a unique subgroup of index \(3\), and \(L/K\) corresponds to this subgroup by Theorem~\ref{thm:conductor}.
\end{proof}

\begin{remark}\label{rem:ray-class-explicit}
The proof gives a concrete presentation of the ray class group:
\[
\Cl_{\mathfrak p_3^2}(K)=\langle [\mathfrak a],u \mid [\mathfrak a]^8=u^3=1,\ [\mathfrak a]u=u[\mathfrak a]\rangle,
\]
where \(\mathfrak a=(2,\tau-1)\) and \(u\) is any nontrivial class in \((1+\mathfrak p_3)/(1+\mathfrak p_3^2)\). The companion audit file records the powers of \(\mathfrak a\) and the corresponding ray-class arithmetic used in Section~5.
\end{remark}

\begin{corollary}\label{cor:weight-one}
Let \(\rho_c\) be the standard two-dimensional Artin representation attached to the splitting field \(L/\QQ\) of \(c(y)\). Then
\[
N(\rho_c)=999,
\qquad
\det(\rho_c)=\chi_{-111},
\]
and there exists a normalized weight-\(1\) cuspidal newform
\[
f_c(q)=\sum_{n\geq 1} a_n q^n\in S_1(\Gamma_0(999),\chi_{-111})
\]
whose \(L\)-function equals \(L(\rho_c,s)\).
\end{corollary}

\begin{proof}
The Artin conductor is \(|\Disc(F/\QQ)|=999\) by Proposition~\ref{prop:number-field}, and the determinant is the quadratic character attached to \(K=\QQ(\sqrt{-111})\). Writing
\[
\rho_c\cong \operatorname{Ind}_{G_K}^{G_\QQ}\chi
\]
for a finite-order character \(\chi\) of \(\Gal(L/K)\), we see that \(\rho_c\) has finite image. Equivalently, \(\rho_c\) is the nontrivial two-dimensional summand of the permutation representation of \(G_\QQ\) on the three roots of \(c(y)\): since \(c(y)\) is irreducible with Galois group \(S_3\), this summand is irreducible. Its determinant is the sign character of \(S_3\), namely \(\chi_{-111}\), so complex conjugation acts with determinant \(-1\) and \(\rho_c\) is odd. The theorem of Deligne--Serre \cite{DeligneSerre1974WeightOne} therefore attaches a normalized weight-\(1\) cuspidal newform of level \(999\) and nebentypus \(\chi_{-111}\) whose \(L\)-function equals \(L(\rho_c,s)\).
\end{proof}

\begin{theorem}\label{thm:hecke-traces}
Let \(f_c(q)=\sum_{n\geq 1} a_n q^n\) be the normalized newform from Corollary~\ref{cor:weight-one}. Then:
\begin{enumerate}
\item[(i)] For every prime \(p\nmid 2\cdot 3\cdot 31\cdot 37\),
\[
a_p=\#\{y\in \F_p:\ c(y)=0\}-1.
\]
\item[(ii)] For every prime \(p\nmid 2\cdot 3\cdot 37\),
\[
a_p=\#\{u\in \F_p:\ u^3-9u^2+256=0\}-1.
\]
\end{enumerate}
In particular, the prime \(31\) is the unique prime away from the Artin conductor at which the maximal-order model \(u^3-9u^2+256\) still detects the Frobenius trace, while the original branch cubic \(c(y)\) does not.
\end{theorem}

\begin{proof}
For \(p\) outside the exceptional set in~(i), the reduction of \(c(y)\) modulo \(p\) is separable, so its factorization type encodes the cycle type of \(\mathrm{Frob}_p\). Let \(\rho_{\mathrm{perm}}\) be the \(3\)-dimensional permutation representation on the roots of~\(c(y)\). Then
\[
\rho_{\mathrm{perm}}\cong \mathbf{1}\oplus \rho_c,
\]
and the trace of \(\rho_{\mathrm{perm}}(\mathrm{Frob}_p)\) is the number of roots of \(c(y)\) fixed by Frobenius, i.e.\ the number of roots of \(c(y)\) in~\(\F_p\). Therefore
\[
a_p=\operatorname{tr}\rho_c(\mathrm{Frob}_p)=\operatorname{tr}\rho_{\mathrm{perm}}(\mathrm{Frob}_p)-1
=\#\mathrm{Fix}(\mathrm{Frob}_p)-1
=\#\{y\in \F_p:\ c(y)=0\}-1.
\]
The same reasoning applies to the monic model \(u^3-9u^2+256\) in~(ii). The exceptional sets differ by the single prime \(31\): the nonmonic model \(c(y)\) drops degree modulo~\(31\) (since \(256\equiv 0\)), while the maximal-order model remains separable there.
\end{proof}
